% !PRINTED 562
\weberpage{562}
%<<CONT
\weq{9}{\begin{aligned}
\eta &= r + r^{-5} + r^{-1} + r^5 = \eta^{(1)}\\
\eta_1 &= r^2 + r^3 + r^{-2} + r^{-3} = \eta^{(2)}\\
\eta_2 &= r^4 + r^6 + r^{-4} + r^{-6} = \eta^{(4)}
\end{aligned}}

Wendet man die Formel (4) an, so erhält man z.\,B.

\[
\begin{aligned}
\eta\eta &= \eta^{(2)} + \eta^{(-4)} + \eta^{(0)} + \eta^{(6)} = 4 + \eta^{(2)} + 2\eta^{(4)}\\
&= -4\eta - 3\eta_1 - 2\eta_2,
\end{aligned}
\]

und so ergeben sich die Formeln

\[
\begin{aligned}
\eta^2 &= -4\eta - 3\eta_1 - 2\eta_2,\\
\eta\eta_1 &= \eta + 2\eta_1 + \eta_2,\\
\eta\eta_2 &= 2\eta + \eta_1 + \eta_2,
\end{aligned}
\]

woraus nach \S.\,167, (17) die Gleichung für $\eta$

\[
\begin{vmatrix}
-4-\eta, & -3, & -2\\
1, & 2-\eta, & 1\\
2, & 1, & 1-\eta
\end{vmatrix} = 0,
\]

oder

\weq{10}{\eta^3 + \eta^2 - 4\eta + 1 = 0}

folgt. Die Discriminante dieser Gleichung ergiebt sich gleich
$169 = 13^2$, also positiv. Die Gleichung hat folglich drei reelle,
% [sic] The print attributes the sign pattern to the positive constant term alone. That
% term shows only that an odd number of roots is negative; the negative coefficient of
% eta supplies the additional exclusion of three negative roots. The conclusion is correct.
und zwar, da das letzte Glied positiv ist, zwei positive und eine
negative Wurzel.

Setzen wir $r = e^{\frac{2\pi i}{13}}$, so wird

\[
\begin{aligned}
\eta &= 2\cos\frac{2\pi}{13} + 2\cos\frac{10\pi}{13} = 2\left(\cos\frac{2\pi}{13} - \cos\frac{3\pi}{13}\right),\\
\eta_1 &= 2\cos\frac{4\pi}{13} + 2\cos\frac{6\pi}{13} = 2\left(\cos\frac{4\pi}{13} + \cos\frac{6\pi}{13}\right),\\
\eta_2 &= 2\cos\frac{8\pi}{13} + 2\cos\frac{12\pi}{13} = -2\left(\cos\frac{\pi}{13} + \cos\frac{5\pi}{13}\right),
\end{aligned}
\]

oder auch

\[
\begin{aligned}
\eta &= 4\cos\frac{4\pi}{13}\cos\frac{6\pi}{13}, \qquad \eta_1 = 4\cos\frac{\pi}{13}\cos\frac{5\pi}{13},\\
\eta_2 &= -4\cos\frac{2\pi}{13}\cos\frac{3\pi}{13}.
\end{aligned}
\]

Es ist also $\eta_2$ die negative, $\eta_1$ die grössere, $\eta$ die kleinere der
beiden positiven Wurzeln.

Die Gleichung $4^{\mathrm{ten}}$ Grades, deren Wurzeln $r$, $r^{-1}$, $r^5$, $r^{-5}$
sind, lässt sich nun leicht bilden, wenn man
%CONT>>
